\begin{center}
\begin{minipage}{\linewidth}\centering
\begin{tikzpicture}[x=1.35cm,y=1.35cm,>=stealth,font=\small,line cap=round,line join=round]
\draw[->,gray!75] (-3.3,0)--(3.46,0) node[right=3pt]{$x$};
\draw[->,gray!75] (0,-3.3)--(0,3.46) node[above=3pt]{$y$};
\node at (-.3,.4) {$O$};
\draw[gray,dashed] (-3.3,-3.3)--(3.3,3.3) node[above left=3pt]{$y=x$};

\draw[teal,very thick,domain=-3:1.65,samples=121] plot(\x,{2^\x});
\draw[blue,very thick,domain=-3:1.65,samples=121] plot({2^\x},\x);
\draw[orange!80!black,very thick,dash dot,domain=-1.65:3,samples=121] plot(\x,{2^(-\x)});
\fill[blue] (2,1) circle(2pt) node[below right=5pt]{$(2,1)$};
\fill[orange!80!black] (2,.25) circle(2pt) node[below=15pt]{$(2,1/4)$};

\end{tikzpicture}\par\smallskip
{\small \textcolor{teal}{\rule[.5ex]{1em}{.8pt}\ $f(x)=2^x$}\qquad\textcolor{blue}{\rule[.5ex]{1em}{.8pt}\ $f^{-1}(x)=\log_2x$}\\[3pt]\textcolor{orange!80!black}{$1/f(x)=2^{-x}$（点划线）}}\par
{\small 镜像得到反函数；取函数值的倒数得到另一条曲线。}\par\smallskip
\end{minipage}
\end{center}
